\chapter{Literature Survey}

The intersection of cryptography and machine learning has garnered significant research interest, particularly concerning the privacy and verifiability of deployed models. Zero-Knowledge Proofs (ZKPs) provide a compelling primitive for these challenges.

\section{Zero-Knowledge Proof Systems}
The foundation of modern ZKPs rests on non-interactive arguments of knowledge (zk-SNARKs). Early constructions like Pinocchio and Groth16 required a trusted setup per circuit, which posed significant logistical and security challenges for deploying multiple complex models~\cite{gennaro2013quadratic}. 

The introduction of PLONK (Permutations over Lagrange-bases for Oecumenical Noninteractive arguments of Knowledge) by Gabizon, Williamson, and Ciobotaru in 2019 revolutionized the space by offering a universal and updatable trusted setup~\cite{kates2020plonk}. This means a single structured reference string (SRS) can be used for any circuit up to a certain size bound. Furthermore, PLONK introduced "custom gates" and lookup tables (Plookup), which are mathematically critical for efficiently evaluating non-algebraic functions like the sigmoid activation within a finite field.

\section{Zero-Knowledge Machine Learning (ZK-ML)}
Applying ZKPs to machine learning (ZK-ML) is an active frontier. Most classical ML models, including logistic regression and deep neural networks, rely heavily on floating-point arithmetic. However, cryptographic circuits operate strictly over prime finite fields.

Recent surveys by Lavin~\cite{lavin2024surveyapplicationszeroknowledgeproofs} outline the primary bottleneck in ZK-ML: approximating non-linear activation functions without exploding the circuit constraint size. Works by Feng et al. have proposed using Taylor series expansions or piecewise polynomial approximations~\cite{cryptoeprint:2023/1345}. While effective for shallow networks, these arithmetic methods degrade either accuracy or performance when implemented as polynomial constraints.

An alternative mechanism is the use of cryptographic lookup tables, as facilitated by the \texttt{gnark} framework's \texttt{logderivlookup} implementation. Lookup tables allow the circuit to replace costly arithmetic polynomial evaluations with an $O(1)$ inclusion check against a pre-computed array of answers.

\section{Scalability and Batch Verification}
While theoretical ZK-ML frameworks exist, practical deployment often stalls due to proving time latency. A single inference proof on a standard dataset can take several seconds, which is unacceptable for real-time APIs. Recent advances focus on proof aggregation and parallelized proving. Although recursive SNARK aggregation offers theoretical constant-size proofs for massive data streams~\cite{cryptoeprint:2023/1174}, the engineering overhead and curve compatibility issues (such as the BLS12-377 to BW6-761 cycle) remain an open challenge in current libraries.

Instead, parallel instance evaluation (batch proving) offers an immediate, scalable engineering solution. By utilizing multi-threaded worker pools in systems like Go, the memory-intensive constraint solving phase can be parallelized, significantly amplifying prediction throughput on modern multi-core architectures such as HPC clusters.
